\section{A continuum without cut points}\label{sec:blocks}

The remaining argument will use paths in a graph whose vertices accumulate
on a continuum without cut points.  We already have such a continuum if
the boundary $B$ has no cut points.  Otherwise, the cut-point tree provides
one of its blocks.  We construct the graph on that block and prove the
topological properties needed below.  In particular, none of the
arguments in this section assumes that $B$ is locally connected.

Recall that $X=\Gamma\sqcup B$ is our compact metric space, with metric
$\rho$, and that $B$ is a nondegenerate continuum.
Propositions~\ref{topo:dimension} and~\ref{topo:cech} give
\[
  \dim B\leq 1,
  \qquad
  \check H^1(B;\mathbb F_2)=0.
\]
Let $P$ be the set of cut points of $B$.  Throughout the construction
below, suppose that $P\ne\varnothing$.  By Proposition~\ref{pret:tree},
there is a nontrivial minimal simplicial $\Gamma$-tree $\mathcal T$
spanned by $P$, with finite quotient and bipartition into $P$ and the
remaining vertices.  All its edge stabilizers are infinite.  We will also
use its separation property: if $q\in P$ separates $a,b\in P$ in $B$,
then $q$ belongs to the open tree interval between $a$ and $b$.
Every point of $P$ is a parabolic fixed point.  In particular,
Lemma~\ref{dyn:parabolic} gives
\begin{equation}\label{block:parabolic-comparison}
  \rho(gp,g)\longrightarrow 0
  \quad\text{as }g\longrightarrow\infty\text{ in }\Gamma,
  \qquad p\in P.
\end{equation}
Here, and below, escape to infinity in a discrete group means escape
from every finite subset.

Choose a vertex $r\in\mathcal T\setminus P$, and write
\[
  R=\{p\in P: p\text{ is adjacent to }r\text{ in }\mathcal T\},
  \qquad H=\Gamma(r).
\]
Since $\mathcal T$ is spanned by $P$, a vertex outside $P$ cannot be a
leaf.  Thus $R$ has at least two points.  For distinct $a,b\in R$, their
tree interval is $a-r-b$.  The separation property therefore shows that
no cut point of $B$ separates $a$ and $b$.  Also, $H$ is exactly the
setwise stabilizer of $R$: an element preserving $R$ preserves the
unique midpoint $r$ of the interval between two distinct members of
$R$.

\subsection{A graph on the cut points}

The graph construction in Lemma~\ref{block:graph} adapts Bowditch's
peripheral-splitting construction~\cite[Section~4, Lemma~4.3 and its
proof]{Bowditch2001Peripheral}. We give the graph argument explicitly.
The subsequent near-path argument and continuous retraction are
extensions established here, not an invocation of a retraction theorem
from that source.

\begin{lemma}\label{block:graph}
There is a connected $\Gamma$-graph $K^*$ with vertex set $P$ and
finitely many edge orbits, such that the endpoints of each edge have
distance two in $\mathcal T$.  The graph $K$ induced by $K^*$ on $R$ is
connected and has finitely many $H$-orbits of vertices and edges.

Define $\pi:P\to R$ by letting $\pi(p)$ be the neighbor of $r$ on the
tree interval from $r$ to $p$.  An edge of $K^*$ with endpoints in $R$
is fixed by $\pi$, and every other edge is collapsed to a vertex.
Moreover, the stabilizer in $H$ of every vertex of $K$ is infinite.
\end{lemma}

\begin{proof}
The finite quotient of $\mathcal T$ gives finitely many $\Gamma$-orbits
in $P$.  Choose representatives $p_0,\ldots,p_k$ and a finite symmetric
generating set $S$ of $\Gamma$.  Start with the graph formed by all
translates of the edges from $p_0$ to $sp_0$, for $s\in S$, and from
$p_0$ to $p_i$, for $1\leq i\leq k$.  Omit loops.  The generator edges
connect the orbit of $p_0$, and the other edges connect every orbit to
it, so this graph is connected.  It has finitely many edge orbits.

Replace each edge by the consecutive pairs of $P$-vertices along the
tree interval between its endpoints.  These intervals are finite, and
the replacement is equivariant.  Thus the resulting graph $K^*$ is
still connected, still has finitely many edge orbits, and has the
required distance-two property.

For an edge $ab$ of $K^*$, let $a-s-b$ be its tree interval.  If $s=r$,
then $a,b\in R$ and $\pi$ fixes both endpoints.  If $s\ne r$, then
$a$ and $b$ lie in the same component of
$\mathcal T\setminus\{r\}$, so $\pi(a)=\pi(b)$.  This proves the
projection assertion.  Projecting a path between two members of $R$
and deleting stationary steps gives a path in $K$, proving its
connectedness.

If a group element carries one edge of $K$ to another, it carries their
unique tree midpoints to one another.  Both midpoints are $r$, so that
element lies in $H$.  Consequently, each $\Gamma$-orbit of edges in
$K^*$ meets $K$ in at most one $H$-orbit.  There are therefore finitely
many $H$-edge orbits in $K$.  Since $K$ is connected and has at least two
vertices, every vertex is incident to an edge; hence there are also
finitely many $H$-vertex orbits.  Finally, for $p\in R$, the stabilizer
$H_p$ is the stabilizer of the tree edge $rp$, which is infinite.
\end{proof}

The graph $K$ need not be locally finite, and its stabilizers need not
be finite.  These facts cause no difficulty: our arguments use finite
orbit sets, rather than finite valence or a proper graph action.

\begin{lemma}\label{block:near-path}
If $a_n,b_n\in P$ both tend to $x\in B$, there are paths in $K^*$ from
$a_n$ to $b_n$ all of whose vertices tend uniformly to $x$.

If the projections of the endpoints of a path in $K^*$ are distinct,
then, after stationary steps are removed, its projected path uses
only vertices of the original path.
\end{lemma}

\begin{proof}
Using the finitely many orbit representatives chosen in
Lemma~\ref{block:graph}, write
$a_n=g_np_{i_n}$ and $b_n=h_np_{j_n}$.  We may choose $g_n$ and $h_n$
to escape every finite subset of $\Gamma$: each representative has
infinite stabilizer, so there are infinitely many choices for each
representing element.  Equation~\eqref{block:parabolic-comparison},
applied to this finite list of representatives, gives
$g_n\to x$ and $h_n\to x$ in $X$.

By Lemma~\ref{compact:near-path}, there are Cayley graph paths from
$g_n$ to $h_n$ whose vertices tend uniformly to $x$.  Fix finite
$K^*$-paths from $p_0$ to $sp_0$ for each Cayley generator $s$, and
from $p_0$ to each representative $p_i$.  Replace a Cayley step
$g\longrightarrow gs$ by the $g$-translate of the first appropriate
path.  At the endpoints, use the translates of the paths to $p_{i_n}$
and $p_{j_n}$.  The result is a $K^*$-path from $a_n$ to $b_n$.

Only finitely many points $t\in P$ occur in the fixed path templates.
For each such point, $\rho(gt,g)\to0$ as $g$ escapes, by
\eqref{block:parabolic-comparison}.  The Cayley path vertices escape
finite sets uniformly, because they tend uniformly to the boundary
point $x$.  Hence all vertices of the transferred paths tend uniformly
to $x$ as well.

For the last assertion, Lemma~\ref{block:graph} says that every
nonconstant projected edge was already an edge of the original path.
When the projected endpoints differ, every vertex of the projected
path after stationary steps are removed is incident to one of these
edges.  It therefore occurred in the original path.
\end{proof}

\subsection{Retraction onto a block}

\begin{proposition}\label{block:retraction}
The map $\pi:P\to R$ extends to a continuous retraction
\[
  \mathfrak r:B\longrightarrow L,
  \qquad L=\overline R^{\,B},
\]
which is locally constant on $B\setminus L$.  The space $L$ is a
nondegenerate continuum without cut points, and
\[
  \dim L\leq1,
  \qquad \check H^1(L;\mathbb F_2)=0.
\]
\end{proposition}

\begin{proof}
By the minimality established in Proposition~\ref{dyn:minimal}, the
nonempty invariant set $P$ is dense in $B$.  We first determine the
limit of $\pi(p_n)$ whenever $p_n\in P$ tends to a point of $B$.

Let $x\in B\setminus L$, and choose a neighborhood $O$ of $x$ in $B$
disjoint from $L$.  The map $\pi$ is constant on $P\cap V$ for some
neighborhood $V\subset O$ of $x$.  Otherwise, we could choose
$a_n,b_n\to x$ with $\pi(a_n)\ne\pi(b_n)$.
Lemma~\ref{block:near-path} would join them by paths whose vertices
eventually all lie in $O$.  Their nonconstant projected paths contain
vertices in $R$ that also belong to the original paths.  This
contradicts $O\cap R=\varnothing$.

Now let $x\in L$ and $a_n\in P$ tend to $x$.  Choose $b_n\in R$ tending
to $x$, and join $a_n$ to $b_n$ by the paths supplied by
Lemma~\ref{block:near-path}.  Whenever $\pi(a_n)\ne b_n$, the same
Lemma places $\pi(a_n)$ on the original path.  Therefore
$\pi(a_n)\to x$.  The indices for which $\pi(a_n)=b_n$ satisfy the
same conclusion directly.

Take the closure of the graph of $\pi$ in $B\times L$.  Density of
$P$ and compactness of $L$ show that this closed set projects onto
$B$.  The preceding two paragraphs show that it has exactly one point
over each point of $B$.  Its projection onto $B$ is consequently a
continuous bijection from a compact space to a Hausdorff space, hence
a homeomorphism.  The second coordinate defines the desired
continuous map $\mathfrak r$.  It is the identity on $L$, and the
first paragraph makes it locally constant on $B\setminus L$.

Since $B$ is connected, its image $L$ is connected.  The set $R$ has at
least two points, so $L$ is nondegenerate.  The dimension bound follows
because $L$ is a closed subset of the compact metric space $B$.
If $i:L\hookrightarrow B$ denotes inclusion, then
$\mathfrak r\circ i=\operatorname{id}_L$.  Functoriality of \v Cech
cohomology gives
$i^*\mathfrak r^*=\operatorname{id}_{\check H^1(L;\mathbb F_2)}$.
Thus $\mathfrak r^*$ injects this group into
$\check H^1(B;\mathbb F_2)=0$.

It remains to rule out a cut point of $L$.  Suppose that
$L\setminus\{q\}=U\sqcup V$ is a separation into nonempty open sets.
Density supplies $a\in R\cap U$ and $b\in R\cap V$.  The extra fiber
\[
  F=\mathfrak r^{-1}(\{q\})\setminus\{q\}
\]
is closed in $B\setminus\{q\}$ by continuity.  It is also open there:
it lies in $B\setminus L$, where $\mathfrak r$ is locally constant.
Consequently,
\[
 B\setminus\{q\}
   =\bigl(\mathfrak r^{-1}(U)\cup F\bigr)
      \sqcup\mathfrak r^{-1}(V)
\]
is a separation, with $a$ and $b$ on opposite sides.  This would make
$q$ a cut point of $B$ separating two members of $R$, contrary to the
separation property of the tree noted above.  Hence $L$ has no cut
points.
\end{proof}

\subsection{The common setting for path pushing}

We now include the case $P=\varnothing$ and collect the exact inputs
for path pushing. The first two properties below retain the boundary
topology and convergence dynamics. The graph conditions describe how
vertices and edges approach the boundary: finite orbit sets and comparison
with group vertices will control edge lengths even when the graph has
infinite valence. When graph vertices lie in the boundary, stabilizer
powers send any prescribed finite set of boundary points toward the
vertex they fix; this will connect
the convergence dynamics to graph components during path pushing.
All closures and metrics in the following proposition are inherited
from $X$.

\begin{proposition}\label{block:setting}
There are a subgroup $H\leq\Gamma$, a compact space
$Y=H\sqcup L\subset X$, and a connected $H$-graph $K$ with vertex set
$I\subset Y$, with the following properties.
\begin{enumerate}
\item $Y=\overline H^{\,X}$.  The points of $H$ are isolated in $Y$,
and $L$ is a nondegenerate continuum without cut points satisfying
\[
  \dim L\leq1,
  \qquad \check H^1(L;\mathbb F_2)=0.
\]
\item If $h_n\to a\in L$ and $h_n^{-1}\to b\in L$, then $h_n$
converges to the constant map $a$, uniformly on compact subsets of
$Y\setminus\{b\}$.
\item The graph $K$ has finitely many $H$-orbits of vertices and
edges.  Either $I=H$, or $I\subset L$.  Every point of $L$ is the limit
of a sequence of distinct vertices in $I$.
\item For each fixed $w\in I$,
\begin{equation}\label{block:comparison}
  \rho(hw,h)\longrightarrow0
  \quad\text{as }h\longrightarrow\infty\text{ in }H.
\end{equation}
\item For every $x\in L$ and every neighborhood $W$ of $x$ in $Y$,
there is a neighborhood $V$ of $x$ in $Y$ such that any two vertices
of $I\cap V$ can be joined by a path in $K$ all of whose vertices
belong to $W$.
\item If $I\subset L$, then for each $p\in I$ there is an
infinite-order element $h\in H_p$ such that
\[
  h^jz\longrightarrow p
  \quad(j\longrightarrow+\infty),\qquad z\in L.
\]
\end{enumerate}
No finite generation assumption on $H$ or local finiteness assumption
on $K$ is included in these conclusions.
\end{proposition}

\begin{proof}
If $P=\varnothing$, take $H=\Gamma$, $Y=X$, $L=B$, and let $K$ be the
original Cayley graph, with $I=H$.  The continuum properties are those
already proved for $B$, and absence of cut points is the definition of
this case.  The closure, convergence, finite orbit, and near-path
assertions follow from the construction of $X$,
Proposition~\ref{dyn:convergence}, and
Lemma~\ref{compact:near-path}.  Every boundary point is approached by
distinct group vertices, since the group vertices are isolated and
dense in $X$.

For completeness, the comparison assertion in this case also follows
from convergence.  If it failed for a fixed $w\in H$, choose escaping
$h_n$ with $\rho(h_nw,h_n)$ bounded away from zero.  After passing to a
subsequence, compactness gives $h_n\to a\in B$ and
$h_n^{-1}\to b\in B$.  Since $w\in H$ is different from $b$,
Proposition~\ref{dyn:convergence} gives $h_nw\to a$, a contradiction.
The last assertion has no content in this case.

If $P\ne\varnothing$, use the objects constructed above:
\[
  H=\Gamma(r),\qquad I=R,\qquad
  L=\overline R^{\,B},\qquad Y=H\sqcup L,
\]
and let $K$ be the graph of Lemma~\ref{block:graph}.
The continuum and finite orbit assertions follow from that Lemma and
Proposition~\ref{block:retraction}.  We verify the remaining
interfaces explicitly.

First, $Y=\overline H^{\,X}$.  Fix $p\in R$.  If an escaping sequence
$h_n\in H$ converges to a boundary point, then
\eqref{block:parabolic-comparison} and $h_np\in R$ show that its limit
belongs to $L$.  Conversely, for any $p\in R$, the infinite stabilizer
$H_p$ supplies distinct elements $h_n$ with $h_np=p$.
Equation~\eqref{block:parabolic-comparison} gives $h_n\to p$.  Thus
$R\subset\overline H^{\,X}$, and taking closures gives
$L\subset\overline H^{\,X}$.  No point of $\Gamma\setminus H$ belongs
to this closure, since all group vertices are isolated.  This proves
the equality and compactness of $Y$.  The subgroup $H$ preserves both
$L$ and $Y$, and the convergence assertion is the restriction of
Proposition~\ref{dyn:convergence}.

The set $R$ is dense in the nondegenerate continuum $L$, which has no
isolated points.  It follows that every point of $L$ is approached by
distinct members of $R$: otherwise a sufficiently small neighborhood
would meet $R$ in at most one point and make that point isolated in
$L$.  Equation~\eqref{block:comparison} is exactly
\eqref{block:parabolic-comparison} restricted to $H$.

For the near-path assertion, take two sequences in $R$ tending to the
same point of $L$.  Lemma~\ref{block:near-path} joins them by $K^*$-paths
whose vertices tend uniformly to that point.  Project to $R$.  If the
endpoints differ, the projected paths use only vertices of the
original paths; if they agree, use the constant path.  Thus the
projected $K$-paths still tend uniformly to that point.  In metric
spaces this sequential statement gives the stated neighborhood
formulation: a failure for some $W$ would yield endpoints in
successively smaller balls that cannot be joined inside $W$.

Finally, $H_p$ is infinite by Lemma~\ref{block:graph}.  Since $\Gamma$
is torsion-free, it contains an infinite-order element.  The
cut-point stabilizer conclusions in Lemma~\ref{pret:parabolic-cuts}
make this element parabolic with fixed point $p$.  Its powers have
the asserted dynamics by Lemma~\ref{dyn:parabolic}.
\end{proof}
